\section{Introducción: no es una entrevista más sobre una startup de generación de video}

Esta sección explica por qué vale la pena convertir EP123 en un material sobre economía de plataformas y métodos de producto de IA, en lugar de limitarse a registrar la historia emprendedora de ONE2X y Mideo. El hilo conductor de Wang Guan (王冠) tiene, de hecho, tres niveles: primero, cómo cambia el papel del gerente de producto de IA en la era de los modelos grandes; segundo, por qué los datos determinan los límites de la inteligencia y cómo las empresas de aplicaciones encuentran oportunidades mediante los datos que genera el propio producto; tercero, cómo los sistemas generativos reescribirán los sistemas de recomendación, las plataformas de contenido y la distribución del poder entre creadores, plataformas y consumidores.

La frase más contundente del episodio es «sin intermediarios que se queden con el margen». Aquí el intermediario no es una empresa concreta, sino el papel de asignador de inventario que las plataformas de internet han desempeñado durante mucho tiempo: los creadores producen el contenido y la plataforma se encarga de distribuirlo, emparejarlo, ordenarlo, cobrar comisión y monetizarlo. El impacto de los sistemas generativos consiste en que interiorizan la parte de «asignación» dentro de la «producción bajo demanda»: el usuario ya no necesariamente entra primero al acervo de inventario a buscar contenido, sino que entrega su intención al sistema de producción, y este genera directamente contenido que corresponde a lo que el usuario necesita.

\begin{importantbox}{Tesis central del episodio}
En la era de la IA, los bienes de información pasarán de una «economía de plataformas de distribución» a una «economía de sistemas de producción». El contenido dejará de ser solo una obra aislada que entra al fondo de tráfico y podrá convertirse en una IP continua, generada en conjunto por la recipe, el context, el modelo y la intención del usuario; la moneda de las plataformas también podría pasar de la atención a la confianza.
\end{importantbox}

\begin{knowledgebox}{Nota sobre la estrategia visual}
Este video es una entrevista con toma fija, sin slides, pizarrón ni demostración de producto. El texto usa la portada para identificar la fuente, y todas las imágenes del cuerpo son diagramas conceptuales de elaboración propia, que sirven para explicar mecanismos abstractos como las tres etapas de los datos, la pila de un sistema generativo, las diferencias entre sistemas de recomendación y sistemas generativos, y el desplazamiento del poder de las plataformas.
\end{knowledgebox}

\subsection{Resumen de la sección}

El valor de EP123 está en discutir de manera conjunta los productos de IA, el ciclo virtuoso de datos y la estructura de poder de las plataformas de contenido. No se pregunta solo «cómo se hace la generación de video», sino: cuando los sistemas generativos asuman la producción de contenido, ¿cómo cambiarán los papeles de las plataformas, los creadores, los consumidores y los gerentes de producto?
